\clearpage\thispagestyle{empty}
\pdfbookmark[0]{Quellen zur Besprechung der Mailänder Ephemeriden für 1816}{carlini1816-sources}
\begin{center}\Large\bfseries Quellen zur Besprechung der Mailänder Ephemeriden für 1816\end{center}

Die Übertragung folgt Carl Friedrich Gauß’ Besprechung in den \emph{Werken}, Bd.~VI (Göttingen, 1874), S.~592--593. In der Werkausgabe lautet die Seitenangabe \emph{457..46c} mit einem offenen, c-förmigen Schlusszeichen. Der Erstabdruck erschien in den \emph{Göttingischen gelehrten Anzeigen}, Stück~47, 23.~März 1816, S.~457--460.

Bei der Angabe zur Libelle stehen im Erstabdruck (S.~459) $1\frac{2}{3}$ Linien. Der Nenner in der Werkausgabe ist undeutlich. Die Übertragung gibt seinen Umriss wieder. Der Dezimalpunkt in der Exzentrizität der Vesta und die offene Klammer bei der Bahnneigung folgen ebenfalls der Werkausgabe. Die historischen Maßeinheiten bleiben erhalten. Cesaris gibt die Regenmengen in Pariser Zoll und Linien an (Anhang S.~88). Die Formen \emph{nordliche}, \emph{Erfahrnngen}, \emph{Mitagsfernrohr} und \emph{häufigern} werden in moderner Schreibweise wiedergegeben.

Das Titelblatt von Francesco Carlinis \emph{Effemeridi astronomiche di Milano per l’anno bisestile 1816} nennt Mailand, Cesarea Regia Stamperia, 1815. Das von Google digitalisierte und über das INAF-Observatorium Brera zugängliche Exemplar trägt einen Stempel der Königlichen Universität Rom. Sein Inhaltsverzeichnis nennt die drei Zusätze zum Kalender und die Beiträge des Anhangs. Orianis erster Aufsatz beschreibt einen Repetitionskreis mit drei Fuß Durchmesser. Die Zusammenfassung seiner Sonnenwendbeobachtungen steht auf S.~85 des Anhangs.

KI-Übertragung aus dem historischen Deutsch in modernes Deutsch. KI-erstellter Satz.

\href{https://gdz.sub.uni-goettingen.de/id/PPN236007467}{Werkausgabe}. Erstdruck, S.~\href{https://images.p.sub.uni-goettingen.de/iiif/image/gdz:PPN319721507_1816:00000459/full/1650,/0/default.jpg}{457}, \href{https://images.p.sub.uni-goettingen.de/iiif/image/gdz:PPN319721507_1816:00000460/full/1650,/0/default.jpg}{458}, \href{https://images.p.sub.uni-goettingen.de/iiif/image/gdz:PPN319721507_1816:00000461/full/1650,/0/default.jpg}{459}, \href{https://images.p.sub.uni-goettingen.de/iiif/image/gdz:PPN319721507_1816:00000462/full/1650,/0/default.jpg}{460}. \href{https://www1.brera.inaf.it/~mario.carpino/Biblioteca_Digitale/Effemeridi_Milano/mirror/EM1816b.pdf#page=2}{Titelblatt des Jahrgangs}. \href{https://www1.brera.inaf.it/~mario.carpino/Biblioteca_Digitale/Effemeridi_Milano/mirror/EM1816b.pdf#page=8}{Inhaltsverzeichnis}. \href{https://www1.brera.inaf.it/~mario.carpino/Biblioteca_Digitale/Effemeridi_Milano/mirror/EM1816b.pdf#page=120}{Oriani zur Refraktion}. \href{https://www1.brera.inaf.it/~mario.carpino/Biblioteca_Digitale/Effemeridi_Milano/mirror/EM1816b.pdf#page=130}{Cesaris zur Bewegung der Gebäude}. \href{https://www1.brera.inaf.it/~mario.carpino/Biblioteca_Digitale/Effemeridi_Milano/mirror/EM1816b.pdf#page=150}{Carlinis Vesta-Tafel}. \href{https://www1.brera.inaf.it/~mario.carpino/Biblioteca_Digitale/Effemeridi_Milano/mirror/EM1816b.pdf#page=202}{Orianis Ergebnisse der Sonnenwendbeobachtungen}. \href{https://www1.brera.inaf.it/~mario.carpino/Biblioteca_Digitale/Effemeridi_Milano/mirror/EM1816b.pdf#page=205}{Cesaris zum Niederschlag}.

\clearpage
